% !PRINTED 77
\weberpage{77}
%<<CONT
\weq{2}{\vartheta_{01}^2\,\vartheta_{01}(u+v)\,\vartheta_{01}(u-v) = \vartheta_{01}^2(u)\,\vartheta_{01}^2(v) - \vartheta_{11}^2(u)\,\vartheta_{11}^2(v).}
\weq{3}{\vartheta_{10}^2\,\vartheta_{10}(u+v)\,\vartheta_{10}(u-v) = \vartheta_{10}^2(u)\,\vartheta_{10}^2(v) - \vartheta_{11}^2(u)\,\vartheta_{11}^2(v).}
\weq{4}{\vartheta_{01}^2\,\vartheta_{11}(u+v)\,\vartheta_{11}(u-v) = \vartheta_{11}^2(u)\,\vartheta_{01}^2(v) - \vartheta_{01}^2(u)\,\vartheta_{11}^2(v).}

These four equations can, as has already been mentioned, be transformed in manifold ways by applying the formulae (13) of the preceding section. The six formulae that follow have only \gesperrt{one} form. Beginning again with an arbitrary example, we have:
\[
\vartheta_{00}(u+v)\,\vartheta_{11}(u-v) = A\,\vartheta_{01}(v)\,\vartheta_{10}(v) + B\,\vartheta_{00}(v)\,\vartheta_{11}(v),
\]
and if one sets $v=0$:
\[
A\,\vartheta_{01}\vartheta_{10} = \vartheta_{00}(u)\,\vartheta_{11}(u),
\]
from which one obtains $B$ by interchanging $u$ with $v$. In this way the three following formulae are derived:

\weq{5}{\begin{aligned}\vartheta_{01}\vartheta_{10}\,\vartheta_{00}(u+v)\,\vartheta_{11}(u-v)\\ {}= \vartheta_{00}(u)\vartheta_{11}(u)\vartheta_{01}(v)\vartheta_{10}(v) - \vartheta_{01}(u)\vartheta_{10}(u)\vartheta_{00}(v)\vartheta_{11}(v),\end{aligned}}
\weq{6}{\begin{aligned}\vartheta_{00}\vartheta_{01}\,\vartheta_{10}(u+v)\,\vartheta_{11}(u-v)\\ {}= \vartheta_{10}(u)\vartheta_{11}(u)\vartheta_{00}(v)\vartheta_{01}(v) - \vartheta_{00}(u)\vartheta_{01}(u)\vartheta_{10}(v)\vartheta_{11}(v),\end{aligned}}
\weq{7}{\begin{aligned}\vartheta_{00}\vartheta_{10}\,\vartheta_{01}(u+v)\,\vartheta_{11}(u-v)\\ {}= \vartheta_{01}(u)\vartheta_{11}(u)\vartheta_{00}(v)\vartheta_{10}(v) - \vartheta_{00}(u)\vartheta_{10}(u)\vartheta_{01}(v)\vartheta_{11}(v),\end{aligned}}
which can be verified immediately by setting first $u$, then $v=0$. From these the three others follow by increasing $u$ by
\[
\frac{1}{2}, \quad \frac{\omega}{2}, \quad \frac{1}{2}:
\]
% [sic] German source repeats 1/2 as the third shift.
\weq{8}{\begin{aligned}\vartheta_{01}\vartheta_{10}\,\vartheta_{01}(u+v)\,\vartheta_{10}(u-v)\\ {}= \vartheta_{01}(u)\vartheta_{10}(u)\vartheta_{01}(v)\vartheta_{10}(v) + \vartheta_{00}(u)\vartheta_{11}(u)\vartheta_{00}(v)\vartheta_{11}(v),\end{aligned}}
\weq{9}{\begin{aligned}\vartheta_{00}\vartheta_{01}\,\vartheta_{00}(u+v)\,\vartheta_{01}(u-v)\\ {}= \vartheta_{00}(u)\vartheta_{01}(u)\vartheta_{00}(v)\vartheta_{01}(v) - \vartheta_{10}(u)\vartheta_{11}(u)\vartheta_{10}(v)\vartheta_{11}(v),\end{aligned}}
\weq{10}{\begin{aligned}\vartheta_{00}\vartheta_{10}\,\vartheta_{00}(u+v)\,\vartheta_{10}(u-v)\\ {}= \vartheta_{00}(u)\vartheta_{10}(u)\vartheta_{00}(v)\vartheta_{10}(v) + \vartheta_{01}(u)\vartheta_{11}(u)\vartheta_{01}(v)\vartheta_{11}(v).\end{aligned}}

These formulae, whose totality is designated by the name of the addition theorem, admit of generalisation in manifold ways, of which the following may serve as an example.

If $u$, $v$ are two variables, then the four products
\[
\begin{array}{ll}
\vartheta_{00}(u)\,\vartheta_{00}(u+v), & \vartheta_{01}(u)\,\vartheta_{01}(u+v), \\
\vartheta_{10}(u)\,\vartheta_{10}(u+v), & \vartheta_{11}(u)\,\vartheta_{11}(u+v),
\end{array}
\]
%CONT>>
